% Breakdown and leakage (bjt-breakdown): the set-up a person would wire, semantically the deck's.
%   RBA ba 0 1e6     the base tied to the emitter through 1 MΩ                     -> resistor R_BA
%   VA sw ca 0       a 0 V source from the swept node to the collector: reads I   -> ammeter I (measured)
%   VSW sw 0 0, .dc  the collector supply, swept                                   -> source V_CE (set)
%   QA ca ba 0       the transistor, emitter grounded                              -> npn
%   The collector-emitter deck is drawn; collector-base and emitter-base are two more.
\documentclass[tikz, border=10pt]{standalone}
\usepackage{bench}
\begin{document}
\begin{tikzpicture}
\benchtitle{0,7.0}{Breakdown and leakage}
\def\gnd{1.0}
% the transistor
\node[npn] (Q) at (4,3) {};
\node[anchor=west, font=\small, align=left] at ($(Q.C)+(0.15,-0.55)$) {DUT};
% emitter to ground
\draw (Q.E) -- (Q.E |- 0,\gnd) node[ground]{};
\nodename{anchor=west}{$(Q.E)+(0.1,-0.3)$}{0}
% the ammeter in the collector, up to the rail
\coordinate (ctop) at ($(Q.C)+(0,1.7)$);
\ammeter{Q.C}{ctop}{$I$}
\nodename{anchor=west}{$(Q.C)+(0.1,0.25)$}{ca}
% the supply at the right, + up, - to ground
\coordinate (vtop) at ($(ctop)+(2.6,0)$);
\draw (ctop) -- (vtop);
\draw[set] (vtop) to[V, l=$V_{CE}$] (vtop |- 0,\gnd) node[ground, black]{};
\nodename{anchor=south}{$(ctop)+(1.3,0.05)$}{sw}
% the base tied to the emitter through a resistor
\coordinate (b) at (Q.B);
\coordinate (bsrc) at ($(b)+(-1.4,0)$);
\draw (b) -- (bsrc);
\draw (bsrc) to[R, l_=$R_{BA}$] (bsrc |- 0,\gnd) node[ground]{};
\nodename{anchor=south}{$(b)+(-0.5,0.02)$}{ba}
% the definitions
\definitions{0,-0.1}{%
  $V_{CE}$ swept 0 to 300 V in 1 V steps;\quad $R_{BA} = 1$ M\ohm\\[2pt]
  read $I$;\quad $V_{(BR)CEO}$ = the $V_{CE}$ at which $I$ passes 1 mA, else absent;\quad $I_{CEO} = I$ at 10 V\\[2pt]
  collector-base (emitter open) and emitter-base (collector open): two more decks%ENDDEFS
}
\end{tikzpicture}
\end{document}
